% chapter1banner.tex
% Standalone chapter-opening page: badge + title + photo + objectives + intro.
% This page is entirely separate from chapter1.tex's lesson content.

\khmerornatetitle{ជំពូកទី១៖ ទែម៉ូឌីណាមិច}{១. ទែម៉ូឌីណាមិក}
\khmerbannerphoto{img/1stchapterbanner.jpg}

% --- CHAPTER LIST (2-column grid, edit freely) ---
\begin{center}
\begin{tabularx}{\textwidth}{@{} >{\raggedright\arraybackslash}X >{\raggedright\arraybackslash}X @{}}
  \khmercircnum{1}~ ទ្រឹស្តីស៊ីនេទិកនៃឧស្ម័ន
  &
  \khmercircnum{2}~ ច្បាប់ទី១ទែម៉ូឌីណាមិក
  \\[3pt]
  \khmercircnum{3}~ ម៉ាសុីន
  &
\end{tabularx}
\end{center}

\vspace{0.15cm}

% --- Intro paragraphs ---
\raggedright
\setlength{\parindent}{1.2em}
សូមស្វាគមន៍មកកាន់ថ្នាក់ទី១២! ជំពួកទី១ នៃមុខវិជ្ជារូបវិទ្យាជាជំពួកដែលសំខាន់ដែលយើងនឹងសិក្សាពីពំនោលនៃឧស្ម័នបរិសុទ្ធ និងលក្ខខណ្ឌផ្សេងៗដែលទាក់ទងនិងឧស្ម័នបរិសុទ្ធ ដូចជាល្បឿនប្រសិទ្ធិនៃឧស្ម័នបរិសុទ្ធ សម្ពាធនៃឧស្ម័នបរិសុទ្ធ និងជាពិសេសគឺសមីការភាពនៃឧស្ម័នបរិសុទ្ធដែលបានសិក្សាកាលពីថ្នាក់ទី១១។ បន្ថែមពីលើទ្រឹស្តីសុីនេទិចនៃឧស្ម័ន យើងនឹងសិក្សាបន្ថែមពីលើថ្នាក់ទី១១ ដែលមានច្បាប់ទី១ទែម៉ូឌីណាមិក និងម៉ាសុីនដែលនឹងបង្ហាញពីរបៀបក្នុងការបម្លែងថាមពលកម្តៅឲ្យទៅជាកម្មន្តនិងថាមពលផ្សេងៗក្នុងជីវិតរស់នៅប្រចាំថ្ងៃ។

\clearpage